% !TeX program = xelatex
% UTF-8; compile twice with XeLaTeX. No external images or data files needed.
\documentclass[a4paper,landscape,10pt]{article}
\usepackage[margin=14mm,headheight=14pt,headsep=5mm,footskip=7mm]{geometry}
\usepackage{fontspec}
\setmainfont{DejaVu Sans}
\newfontfamily\ipafont{DejaVu Sans}
\usepackage[table]{xcolor}
\usepackage{array,longtable,booktabs,fancyhdr}
\definecolor{ink}{HTML}{18384D}
\definecolor{band}{HTML}{E7EFF4}
\pagestyle{fancy}\fancyhf{}
\fancyhead[L]{\small\color{ink}Parlez Vous QSO?! · Srpski / Serbian}
\fancyhead[R]{\small Englische Phrase · Serbische Übersetzung · IPA}
\fancyfoot[L]{\scriptsize Vorlage: Serbian v 1.0 · Credits: Marko Milenković -- YT8CDC}
\fancyfoot[R]{\small\thepage}
\setlength{\parindent}{0pt}
\setlength{\tabcolsep}{3mm}
\renewcommand{\arraystretch}{1.25}
\newcommand{\ipa}[1]{{\ipafont[#1]}}
\begin{document}
{\LARGE\bfseries\color{ink}Parlez Vous QSO?!}\hfill{\large Srpski / Serbian}

\vspace{2mm}
{\small QSO-Sprachführer mit serbischer IPA-Lautschrift. Die Lautschrift bezieht sich auf die serbische Spalte.}

\vspace{2mm}
{\footnotesize \textbf{Leseschlüssel:} \ipa{ʃ} = š, \ipa{ʒ} = ž, \ipa{ts} = c, \ipa{tʃ} = č, \ipa{tɕ} = ć,
\ipa{dʑ} = đ, \ipa{ɲ} = nj, \ipa{ʎ} = lj, \ipa{x} = h; \ipa{r̩} ist silbisches r.
Vereinfachte Lautfolge für ekavisches Standardserbisch; Betonung, Tonakzente und Vokallängen sind nicht markiert.
\ipa{…} steht für einen Platzhalter.}

\vspace{1mm}
\fontsize{9}{11.8}\selectfont
\begin{longtable}{@{}>{\raggedright\arraybackslash}p{13mm}>{\raggedright\arraybackslash}p{73mm}>{\raggedright\arraybackslash}p{75mm}>{\raggedright\arraybackslash}p{90mm}@{}}
\toprule
\textbf{Ref.} & \textbf{English phrase} & \textbf{Srpski} & \textbf{Lautschrift (IPA)}\\
\midrule
\endfirsthead
\toprule
\textbf{Ref.} & \textbf{English phrase} & \textbf{Srpski} & \textbf{Lautschrift (IPA)}\\
\midrule
\endhead
\midrule
\multicolumn{4}{r}{\scriptsize Fortsetzung auf der nächsten Seite}\\
\endfoot
\bottomrule
\endlastfoot
\rowcolor{band} 1.0 & \textbf{Opening Greetings} & \multicolumn{2}{l}{\textbf{Pozdrav na početku veze}} \\*
1.1 & I am very pleased to meet you. & Drago mi je da se upoznajemo. & \ipa{draɡo mi je da se upoznajemo} \\
1.2 & Thank you for your call & Hvala na pozivu. & \ipa{xvala na pozivu} \\
1.3 & Good morning & Dobro jutro. & \ipa{dobro jutro} \\
1.31 & Good day & Dobar dan. & \ipa{dobar dan} \\
1.32 & Good afternoon & Dobro popodne. & \ipa{dobro popodne} \\
1.33 & Good evening & Dobro veče. & \ipa{dobro vetʃe} \\
1.4 & I know only a little Serbian but I enjoy QSOs in Serbian. & Znam samo malo srpskog, ali volim da radim veze na srpskom jeziku. & \ipa{znam samo malo sr̩pskoɡ ali volim da radim veze na sr̩pskom jeziku} \\
\rowcolor{band} 2.0 & \textbf{Operator} & \multicolumn{2}{l}{\textbf{Operator}} \\*
2.1 & My name is Marko. & Moje ime je Marko. & \ipa{moje ime je marko} \\
2.2 & I have been licensed since 2005. & Posedujem licencu od 2005. & \ipa{posedujem litsentsu od dve xiʎade pete} \\
2.3 & I have been licensed for 5 years. & Licenciran sam već pet godina. & \ipa{litsentsiran sam vetɕ pet ɡodina} \\
2.4 & I live in Sittingbourne. & Živim u Sitingburnu. & \ipa{ʒivim u sitinɡburnu} \\
2.5 & Sittingbourne is a village. & Sitingburn je selo. & \ipa{sitinɡburn je selo} \\
2.51 & Sittingbourne is a town. & Sitingburn je gradić. & \ipa{sitinɡburn je ɡraditɕ} \\
2.52 & Sittingbourne is a city. & Sitingburn je grad. & \ipa{sitinɡburn je ɡrad} \\
2.6 & Sittingbourne is between London and Dover. & Sitingburn se nalazi između Londona i Dovera. & \ipa{sitinɡburn se nalazi izmedʑu londona i dovera} \\
2.7 & I live in the South of England. & Živim na jugu Engleske. & \ipa{ʒivim na juɡu enɡleske} \\
2.71 & I live in the North of England. & Živim na severu Engleske. & \ipa{ʒivim na severu enɡleske} \\
2.72 & I live in the East of England. & Živim na istoku Engleske. & \ipa{ʒivim na istoku enɡleske} \\
2.73 & I live in West of England. & Živim na zapadu Engleske. & \ipa{ʒivim na zapadu enɡleske} \\
2.74 & I live in South East England. & Živim na jugoistoku Engleske. & \ipa{ʒivim na juɡoistoku enɡleske} \\
2.75 & I live in South West England. & Živim na jugozapadu Engleske. & \ipa{ʒivim na juɡozapadu enɡleske} \\
2.76 & I live in North East England. & Živim na severoistoku Engleske. & \ipa{ʒivim na severoistoku enɡleske} \\
2.77 & I live in North West England. & Živim na severozapadu Engleske. & \ipa{ʒivim na severozapadu enɡleske} \\
2.78 & I live in central England. & Živim u centralnoj Engleskoj. & \ipa{ʒivim u tsentralnoj enɡleskoj} \\
2.8 & I live in the country. & Živim na selu. & \ipa{ʒivim na selu} \\
2.81 & I live by the sea. & Živim pored mora. & \ipa{ʒivim pored mora} \\
2.82 & I live near mountains. & Živim u planinskom delu. & \ipa{ʒivim u planinskom delu} \\
2.9 & Sittingbourne is 50Km from London. & Sitingburn je 50 km od Londona. & \ipa{sitinɡburn je pedeset kilometara od londona} \\
2.11 & I am retired. & U penziji sam. & \ipa{u penziji sam} \\
2.12 & My home call sign is M0ECS. & Moj domaći pozivni znak je M0ECS. & \ipa{moj domatɕi pozivni znak je em nula e tse es} \\
\rowcolor{band} 3.0 & \textbf{Station} & \multicolumn{2}{l}{\textbf{Stanica}} \\*
3.1 & My radio is an Icom 706. & Moj primopredajnik je Icom 706. & \ipa{moj primopredajnik je aɪkɒm sedamsto ʃest} \\
3.2 & My station is all homebrew (homemade). & Moj primopredajnik je kućne izrade. & \ipa{moj primopredajnik je kutɕne izrade} \\
3.3 & I am running 5 Watts. & Radim sa 5 W snage. & \ipa{radim sa pet vati snaɡe} \\
3.4 & My antenna is a long wire. & Moja antena je long wire. & \ipa{moja antena je lɒŋ waɪə} \\
3.41 & My antenna is a dipole. & Moja antena je dipol. & \ipa{moja antena je dipol} \\
3.42 & My antenna is a G5RV. & Moja antena je G5RV. & \ipa{moja antena je ɡe pet er ve} \\
3.43 & My antenna is a doublet. & Moja antena je dublet. & \ipa{moja antena je dublet} \\
3.44 & My antenna is a delta loop. & Moja antena je delta lup. & \ipa{moja antena je delta lup} \\
3.45 & My antenna is a vertical. & Moja antena je vertikalka. & \ipa{moja antena je vertikalka} \\
3.46 & My antenna is a log periodic. & Moja antena je log periodik. & \ipa{moja antena je loɡ periodik} \\
3.47 & My antenna is a Yagi beam. & Moja antena je Yagi bim. & \ipa{moja antena je jaɡi bim} \\
3.48 & My antenna is a 5 element beam. & Moja antena je petoelementni bim. & \ipa{moja antena je petoelementni bim} \\
\rowcolor{band} 4.0 & \textbf{Conditions} & \multicolumn{2}{l}{\textbf{Uslovi}} \\*
4.1 & Thank you for my 599 report. & Hvala za 599 raport. & \ipa{xvala za pet devet devet raport} \\
4.2 & The band conditions are good. & Uslovi na bandu su dobri. & \ipa{uslovi na bandu su dobri} \\
4.3 & The band conditions are bad. & Uslovi na bandu su loši. & \ipa{uslovi na bandu su loʃi} \\
4.4 & There is strong interference here. & Imam jaku interferenciju ovde. & \ipa{imam jaku interferentsiju ovde} \\
4.5 & Can you change to the 80m band? & Da li možeš da se prebaciš na 80 m band? & \ipa{da li moʒeʃ da se prebatsiʃ na osamdeset metara band} \\
4.6 & Please change frequency to 3.558 MHz. & Molim te, promeni frekvenciju na 3,558 MHz. & \ipa{molim te promeni frekventsiju na tri zarez pet pet osam meɡaxertsa} \\
\rowcolor{band} 5.0 & \textbf{General Queries and Comments} & \multicolumn{2}{l}{\textbf{Opšte informacije i komentari}} \\*
5.1 & What does <\#\#\#\#\#> mean please? & Šta znači <\#\#\#\#\#>? & \ipa{ʃta znatʃi …} \\
5.2 & Please wait a moment, I need to check my Serbian dictionary. & Molim te sačekaj, moram da proverim moj rečnik srpskog. & \ipa{molim te satʃekaj moram da proverim moj retʃnik sr̩pskoɡ} \\
5.3 & How do you say, <\#\#\#\#\#> in Serbian ? & Kako se kaže <\#\#\#\#\#> na srpskom? & \ipa{kako se kaʒe … na sr̩pskom} \\
5.5 & My favourite modes of operation are CW / RTTY / PSK / MT63 / Olivia / Hellschreiber / Digital Modes / and SSB & Moji omiljeni modovi rada su CW / RTTY / PSK / MT63 / Olivia / Hellschreiber / digitalni modovi i SSB. & \ipa{moji omiʎeni modovi rada su tse duplo ve / er te te ipsilon / pe es ka / em te ʃezdeset tri / olivija / hɛlʃraɪbɐ / diɡitalni modovi i es es be} \\
5.6 & My best DX is Australia. & Moj najbolji DX je Australija. & \ipa{moj najboʎi de iks je australija} \\
5.7 & I enjoy working contests. & Uživam da radim takmičenja. & \ipa{uʒivam da radim takmitʃeɲa} \\
5.8 & I like working QRP. & Volim da radim QRP. & \ipa{volim da radim ku er pe} \\
\rowcolor{band} 6.0 & \textbf{The Weather} & \multicolumn{2}{l}{\textbf{Vremenski uslovi}} \\*
6.1 & It is sunny. & Sunčano je. & \ipa{suntʃano je} \\
6.2 & It is windy. & Vetrovito je. & \ipa{vetrovito je} \\
6.3 & It is raining. & Pada kiša. & \ipa{pada kiʃa} \\
6.4 & It is snowing. & Pada sneg. & \ipa{pada sneɡ} \\
6.5 & It is very hot. & Veoma je vruće. & \ipa{veoma je vrutɕe} \\
6.6 & It is very cold. & Veoma je hladno. & \ipa{veoma je xladno} \\
6.7 & The temperature is 8C. & Temperatura je 8 stepeni (C). & \ipa{temperatura je osam stepeni tselzijusa} \\
6.8 & It is overcast. & Oblačno je. & \ipa{oblatʃno je} \\
6.9 & It is foggy. & Maglovito je. & \ipa{maɡlovito je} \\
6.11 & It has been raining all day. & Pada kiša celog dana. & \ipa{pada kiʃa tseloɡ dana} \\
6.12 & There is 20cm of snow here. & Palo je 20 centimetara snega. & \ipa{palo je dvadeset tsentimetara sneɡa} \\
\rowcolor{band} 8.0 & \textbf{QSL} & \multicolumn{2}{l}{\textbf{QSL}} \\*
8.1 & My QSL card is sure via the bureau. & Moja QSL karta sigurno stiže preko biroa. & \ipa{moja ku es el karta siɡurno stiʒe preko biroa} \\
8.2 & QSL via eQSL & QSL preko e-QSL. & \ipa{ku es el preko i ku es el} \\
8.3 & No eQSL please & Ne primam eQSL. & \ipa{ne primam i ku es el} \\
8.4 & I will send you my card direct. & Poslaću ti QSL direktno. & \ipa{poslatɕu ti ku es el direktno} \\
8.5 & What is your address? & Koja je tvoja adresa? & \ipa{koja je tvoja adresa} \\
8.6 & Please send your QSL card direct. & Molim te, pošalji tvoju QSL direktno. & \ipa{molim te poʃaʎi tvoju ku es el direktno} \\
8.7 & My address is... & Moja adresa je… & \ipa{moja adresa je …} \\
8.8 & My address is shown on QSL.com / EHam.net / ….. & Moja adresa je prikazana na: qsl.com / EHam.net / … & \ipa{moja adresa je prikazana na ku es el tatʃka kom / i hæm tatʃka net / …} \\
8.9 & Sorry I do not send QSL cards. & Žao mi je, ne šaljem QSL. & \ipa{ʒao mi je ne ʃaʎem ku es el} \\
8.11 & Please send QSL card via G1JQH. & Pošalji QSL preko G1JQH. & \ipa{poʃaʎi ku es el preko ɡe jedan jot ku xa} \\
8.12 & My WACRAL / GQRP / DXCC / FISTS number is 1234. & Moj WACRAL / GQRP / DXCC / FISTS broj je 1234. & \ipa{moj wakral / ɡe ku er pe / de iks tse tse / fɪsts broj je xiʎadu dvesta trideset tʃetiri} \\
\rowcolor{band} 9.0 & \textbf{Final Greetings} & \multicolumn{2}{l}{\textbf{Pozdravi}} \\*
9.1 & Thank you. & Hvala. & \ipa{xvala} \\
9.2 & Thank you for a very enjoyable contact. & Hvala na veoma prijatnom kontaktu. & \ipa{xvala na veoma prijatnom kontaktu} \\
9.3 & I enjoyed our chat. & Uživao sam u našem razgovoru. & \ipa{uʒivao sam u naʃem razɡovoru} \\
9.4 & I enjoyed our QSO in Serbian. & Uživao sam u našem QSO na srpskom. & \ipa{uʒivao sam u naʃem ku es o na sr̩pskom} \\
9.5 & It has been great fun. & Bilo je veoma zabavno. & \ipa{bilo je veoma zabavno} \\
9.6 & If you would like a regular contact with me, you can contact me via email. & Ako želiš redovan kontakt sa mnom, možeš me kontaktirati na e-mail. & \ipa{ako ʒeliʃ redovan kontakt sa mnom moʒeʃ me kontaktirati na imeɪl} \\
9.7 & My email address is... & Moja e-mail adresa je... & \ipa{moja imeɪl adresa je …} \\
9.8 & My email address is shown on QSL.com / EHam.net & Moja e-mail adresa je prikazana na qsl.com / EHam.net. & \ipa{moja imeɪl adresa je prikazana na ku es el tatʃka kom / i hæm tatʃka net} \\
9.9 & I look forward to meeting you again. & Unapred se radujem ponovnom susretu. & \ipa{unapred se radujem ponovnom susretu} \\
9.11 & My best wishes to you and your family. & Sve najbolje tebi i tvojoj porodici. & \ipa{sve najboʎe tebi i tvojoj poroditsi} \\
9.12 & Best wishes to your wife. & Sve najbolje tvojoj supruzi. & \ipa{sve najboʎe tvojoj supruzi} \\
9.121 & Best wishes to your husband. & Sve najbolje tvom suprugu. & \ipa{sve najboʎe tvom supruɡu} \\
9.13 & Good luck. & Srećno. & \ipa{sretɕno} \\
9.14 & Good health. & Dobro zdravlje. & \ipa{dobro zdravʎe} \\
9.15 & Goodbye & Doviđenja. & \ipa{dovidʑeɲa} \\
9.16 & I wish you an enjoyable afternoon. & Želim ti prijatno popodne. & \ipa{ʒelim ti prijatno popodne} \\
9.161 & I wish you an enjoyable evening. & Želim ti prijatno veče. & \ipa{ʒelim ti prijatno vetʃe} \\
9.162 & I wish you an enjoyable weekend. & Želim ti prijatan vikend. & \ipa{ʒelim ti prijatan vikend} \\
9.17 & Enjoy your holiday. & Uživaj u praznicima. & \ipa{uʒivaj u praznitsima} \\
9.18 & Happy New Year. & Srećna nova godina. & \ipa{sretɕna nova ɡodina} \\
\rowcolor{band} 10.0 & \textbf{Country} & \multicolumn{2}{l}{\textbf{Države}} \\*
10.1 & Albania & Albanija & \ipa{albanija} \\
10.2 & Austria & Austrija & \ipa{austrija} \\
10.3 & Belgium & Belgija & \ipa{belɡija} \\
10.4 & Belorussia & Belorusija & \ipa{belorusija} \\
10.5 & Bosnia \& Herzegovina & Bosna i Hercegovina & \ipa{bosna i xertseɡovina} \\
10.6 & Bulgaria & Bugarska & \ipa{buɡarska} \\
10.7 & Croatia & Hrvatska & \ipa{xr̩vatska} \\
10.8 & Czech Republic & Češka Republika & \ipa{tʃeʃka republika} \\
10.9 & Denmark & Danska & \ipa{danska} \\
10.11 & Eire & Irska & \ipa{irska} \\
10.12 & England & Engleska & \ipa{enɡleska} \\
10.13 & Estonia & Estonija & \ipa{estonija} \\
10.14 & Finland & Finska & \ipa{finska} \\
10.15 & France & Francuska & \ipa{frantsuska} \\
10.16 & Germany & Nemačka & \ipa{nematʃka} \\
10.17 & Gibraltar & Gibraltar & \ipa{ɡibraltar} \\
10.18 & Greece & Grčka & \ipa{ɡr̩tʃka} \\
10.19 & Holland & Holandija & \ipa{xolandija} \\
10.21 & Hungary & Mađarska & \ipa{madʑarska} \\
10.22 & Iceland & Island & \ipa{island} \\
10.23 & Northern Ireland & Severna Irska & \ipa{severna irska} \\
10.24 & Isle of Man & Ostrvo Man & \ipa{ostr̩vo man} \\
10.25 & Italy & Italija & \ipa{italija} \\
10.26 & Latvia & Letonija & \ipa{letonija} \\
10.27 & Liechtenstein & Lihtenštajn & \ipa{lixtenʃtajn} \\
10.28 & Lithuania & Litvanija & \ipa{litvanija} \\
10.29 & Luxembourg & Luksemburg & \ipa{luksemburɡ} \\
10.31 & Macedonia & Makedonija & \ipa{makedonija} \\
10.32 & Moldavia & Moldavija & \ipa{moldavija} \\
10.33 & Netherlands & Holandija & \ipa{xolandija} \\
10.34 & Norway & Norveška & \ipa{norveʃka} \\
10.35 & Poland & Poljska & \ipa{poʎska} \\
10.36 & Portugal & Portugal & \ipa{portuɡal} \\
10.37 & Romania & Rumunija & \ipa{rumunija} \\
10.38 & Russian Federation & Rusija & \ipa{rusija} \\
10.39 & Scotland & Škotska & \ipa{ʃkotska} \\
10.41 & Slovakia & Slovačka & \ipa{slovatʃka} \\
10.42 & Slovenia & Slovenija & \ipa{slovenija} \\
10.43 & Spain & Španija & \ipa{ʃpanija} \\
10.44 & Svalbard & Svalbard & \ipa{svalbard} \\
10.45 & Sweden & Švedska & \ipa{ʃvedska} \\
10.46 & Switzerland & Švajcarska & \ipa{ʃvajtsarska} \\
10.47 & Turkey & Turska & \ipa{turska} \\
10.48 & Ukraine & Ukrajina & \ipa{ukrajina} \\
10.49 & Wales & Vels & \ipa{vels} \\
\end{longtable}
\normalsize
\section*{Hinweise zur Übertragung und Aussprache}
\fontsize{8.5}{10}\selectfont\raggedright
\textbf{Vorlage und Umfang.} Grundlage ist „Serbian v 1.0.doc“, Credit: Marko Milenkovic, YT8CDC
(hier mit Diakritikum: Milenković). Alle 142 Phrasen und Ländernamen sind enthalten.
Die Nummerierung und ihre Sprünge bleiben erhalten, einschließlich des fehlenden Abschnitts 7.
Die Vorlage verwendet Lateinschrift; diese bleibt erhalten. Serbisch kann ebenso kyrillisch geschrieben werden.
Der Titel steht als „Srpski“ statt „Srpskog“.

\smallskip
\textbf{Schreibungen und Namen.} Fehlende serbische Zeichen č, ć, š, ž und đ wurden wiederhergestellt;
„izmedju“, „Madjarska“ und „Dovidjenja“ stehen als „između“, „Mađarska“ und „Doviđenja“.
Sitingburn/Sitinburg wurde durchgehend zu „Sitingburn“ vereinheitlicht, mit „Sitingburnu“ nach u (2.4).
„delat lup“ ist „delta lup“ (3.44), „inteferenciju“ ist „interferenciju“ (4.4),
„Temeperatura“ ist „Temperatura“ (6.7), „Linhestajn“ ist „Lihtenštajn“ (10.27).
Sprachbezeichnungen stehen klein; Abstände und Satzzeichen wurden vereinheitlicht.
Der Platzhalter in 5.1 ist an die englische Spalte angeglichen.

\smallskip
\textbf{Inhaltliche Korrekturen.} „na kopnu“ (auf dem Festland) wurde gemäß „in the country“
zu „na selu“ (auf dem Land, 2.8). Die doppelte Angabe MT63 in 5.5 wurde entsprechend der englischen Spalte
zu Hellschreiber; „Digitalni mod“ steht im Plural als „digitalni modovi“.
„Moja najbolji“ wurde zu „Moj najbolji“ (5.6); das überzählige „e’“ in 8.12 wurde entfernt.
Der italienische Satz „Mi sono divertito molto“ wurde durch „Bilo je veoma zabavno“ ersetzt (9.5).
„drugi kontakt samnom“ wurde gemäß dem regelmäßigen Kontakt der englischen Spalte zu
„redovan kontakt sa mnom“ (9.6). „Srena Nova Godina“ wurde zu „Srećna nova godina“ (9.18).
Die Frequenz in 4.6 verwendet ein Dezimalkomma: 3,558 MHz.

\smallskip
\textbf{Quellenbesonderheiten.} Name Marko, Lizenzjahr 2005 und Dauer fünf Jahre sind Beispiele der Vorlage.
3.2 beschreibt den selbstgebauten Transceiver statt der gesamten Station.
2.82 nennt einen gebirgigen Landesteil statt ausdrücklich die Nähe zu Bergen.
9.17 wünscht Freude an Feiertagen (praznici) statt allgemein am Urlaub.
Die männlichen Formen licenciran und uživao bleiben erhalten; Sprecherinnen verwenden licencirana/uživala.
Historische Ländernamen wie Belorusija und Makedonija sowie die Adressen bleiben unverändert.

\smallskip
\textbf{Aussprachebasis und Umfang.} Die IPA-Spalte gibt die Lautfolge des ekavischen Standardserbisch
vereinfacht wieder, passend zu sever, sneg und vreme in der Vorlage.
Sie ist ausdrücklich keine vollständige prosodische Transkription: Wortbetonung, die vier traditionellen
Tonakzente, betonte und nachtonige Vokallängen sowie Satzmelodie sind nicht markiert.
Fehlende Akzentzeichen bedeuten daher nicht, dass die Wörter unbetont oder alle Vokale kurz wären.
Diese Vereinfachung gilt auch für fremdsprachige Namen und Fachwörter.

\smallskip
\textbf{Laute.} \ipa{tʃ} (č) und \ipa{tɕ} (ć) sowie \ipa{dʒ} (dž) und \ipa{dʑ} (đ)
sind unterschiedliche Affrikaten; c ist \ipa{ts}. Nj \ipa{ɲ} und lj \ipa{ʎ} sind jeweils ein Laut.
R steht vereinfacht als \ipa{r}; in srpski, Hrvatska und Grčka ist es silbisch \ipa{r̩}.
H entspricht \ipa{x}, etwa in hvala. Die Vokale stehen breit als \ipa{a, e, i, o, u}.
Feine Vokalqualitäten, die variable Qualität von l sowie Assimilationen im zusammenhängenden Satz
werden nicht durchgehend dargestellt.

\smallskip
\textbf{Namen und Kürzel.} Orts- und Ländernamen sind gemäß den serbischen Formen gelesen,
auch Sitingburn, Dovera und Lihtenštajn. Long wire, Icom, Hellschreiber, WACRAL, FISTS, Ham und e-mail
verwenden vereinfachte englische beziehungsweise deutsche Lesarten.
Dublet, delta lup, log periodik und jagi bim sind serbisch gelesen.
Buchstabennamen: Q „ku“, S „es“, L „el“, G „ge“, C „ce“, J „jot“, H „ha“, W „duplo ve“, Y „ipsilon“;
andere Lesarten sind möglich. eQSL/eHam verwenden hier ein englisches \ipa{i} für e.
Webpunkte werden als „tačka“ gesprochen, com/net als Wörter.

\smallskip
\textbf{Zahlen.} „od 2005“ = „od dve hiljade pete“, 706 = „sedamsto šest“,
599 = „pet devet devet“, MT63 mit „šezdeset tri“, 1234 = „hiljadu dvesta trideset četiri“.
3,558 MHz wird als „tri zarez pet pet osam megaherca“ gelesen.
50 km, 5 W, 80 m, 8 °C und 20 cm stehen mit gesprochenen Einheiten;
C in 6.7 wird als „Celzijusa“ aufgelöst. Offene Angaben und Platzhalter erscheinen als \ipa{…}.

\smallskip
\textbf{Fachhinweise.} Zum serbischen Akzentsystem: Robin Karlin,
„Effects of syllable onset on the timing of pitch accent in Belgrade Serbian“,
Proceedings of the Linguistic Society of America 3, 2018, 36:1--15, DOI 10.3765/plsa.v3i1.4324.
Zum Kontrast l/lj: „The contrast and positional/contextual variation in Serbian laterals“,
Journal of the International Phonetic Association, Cambridge Core.
\end{document}
